\begin{table}[htbp]
\centering
\scriptsize
\caption{Экспериментальные и рассчитанные центры полос ${}^{76}$GeH$_4$.}
\label{tab:ge_band_centers_76}
\setlength{\tabcolsep}{3pt}
\begin{tabular}{lcccccccc}
\toprule
Полоса & $N$ & $E_{\textrm{эксп}}$ & $J_{\max}$ & $N_{\textrm{пер}}$ & $N_{\textrm{ур}}$ & $E_{\textrm{расч}}$ & $\delta$ & Ист. \\
 & & \multicolumn{1}{c}{см$^{-1}$} & & & & \multicolumn{1}{c}{см$^{-1}$} & \multicolumn{1}{c}{см$^{-1}$} & \\
\midrule
$\nu_4$($F_2$) & 1/2 & 820.33 & 26 & 1117 & 588 & 820.38 & -0.05 & \cite{871} \\
$\nu_2$($E$) & 1/2 & 929.91 & 23 & 805 & 318 & 929.95 & -0.04 & \cite{871} \\
$2\nu_4$($A_1$) & 1 & 1627.50 & 11 & 76 & 29 & 1627.50 & 0.00 & \cite{872} \\
$2\nu_4$($F_2$) & 1 & 1639.26 & 15 & 306 & 189 & 1639.24 & 0.02 & \cite{872} \\
$2\nu_4$($E$) & 1 & 1642.14 & 13 & 80 & 43 & 1642.22 & -0.08 & \cite{872} \\
$\nu_2+\nu_4$($F_2$) & 1 & 1748.40 & 19 & 1148 & 401 & 1748.48 & -0.08 & \cite{872} \\
$\nu_2+\nu_4$($F_1$) & 1 & 1752.50 & 18 & 923 & 289 & 1752.54 & -0.04 & \cite{872} \\
$2\nu_2$($A_1$) & 1 & 1857.27 & 18 & 131 & 64 & 1857.26 & 0.01 & \cite{872} \\
$2\nu_2$($E$) & 1 & 1860.67 & 19 & 584 & 257 & 1860.68 & -0.01 & \cite{872} \\
$\nu_1$($A_1$) & 1 & 2110.69 & 20 & 445 & 149 & 2110.93 & -0.24 & \cite{873} \\
$\nu_3$($F_2$) & 1 & 2110.73 & 31 & 1366 & 714 & 2110.73 & 0.00 & \cite{873} \\
$3\nu_4$($F_2$) & 3/2 & 2438.24 & 20 & 732 & 441 & 2438.24 & 0.00 & \cite{874} \\
$3\nu_4$($A_1$) & 3/2 & 2456.25$^{(*)}$ &  &  &  & 2456.57 & -0.32 & \cite{874} \\
$3\nu_4$($F_1$) & 3/2 & 2459.51 & 16 & 147 & 91 & 2459.56 & -0.05 & \cite{874} \\
$3\nu_4$($F_2$) & 3/2 & 2462.40 & 17 & 242 & 117 & 2462.31 & 0.09 & \cite{874} \\
$\nu_2+2\nu_4$($E$) & 3/2 & 2556.20 & 14 & 223 & 105 & 2556.77 & -0.57 & \cite{874} \\
$\nu_2+2\nu_4$($F_1$) & 3/2 & 2567.68 & 14 & 192 & 104 & 2567.54 & 0.14 & \cite{874} \\
$\nu_2+2\nu_4$($A_1$) & 3/2 & 2568.97$^{(*)}$ &  &  &  & 2568.48 & 0.49 & \cite{874} \\
$\nu_2+2\nu_4$($F_2$) & 3/2 & 2571.48 & 16 & 137 & 84 & 2571.59 & -0.11 & \cite{874} \\
$\nu_2+2\nu_4$($E$) & 3/2 & 2573.44 & 13 & 43 & 29 & 2573.59 & -0.15 & \cite{874} \\
$\nu_2+2\nu_4$($A_2$) & 3/2 & 2576.54$^{(*)}$ &  &  &  & 2576.60 & -0.06 & \cite{874} \\
$2\nu_2+\nu_4$($F_2$) & 3/2 & 2675.23 & 19 & 564 & 329 & 2675.31 & -0.08 & \cite{874} \\
$2\nu_2+\nu_4$($F_1$) & 3/2 & 2681.39 & 15 & 320 & 167 & 2681.43 & -0.04 & \cite{874} \\
$2\nu_2+\nu_4$($F_2$) & 3/2 & 2684.19 & 16 & 355 & 141 & 2684.13 & 0.06 & \cite{874} \\
$3\nu_2$($E$) & 3/2 & 2785.41$^{(*)}$ &  &  &  & 2785.35 & 0.06 & \cite{874} \\
$3\nu_2$($A_1$) & 3/2 & 2791.95$^{(*)}$ &  &  &  & 2792.18 & -0.23 & \cite{874} \\
$3\nu_2$($A_2$) & 3/2 & 2792.15$^{(*)}$ &  &  &  & 2792.18 & -0.03 & \cite{874} \\
$\nu_3+\nu_4$($F_1$) & 3/2 & 2924.24 & 20 & 937 & 446 & 2924.24 & 0.00 & \cite{875} \\
$\nu_3+\nu_4$($E  $) & 3/2 & 2924.99 & 18 & 124 & 64 & 2924.99 & 0.00 & \cite{875} \\
$\nu_1+\nu_4$($F_2$) & 3/2 & 2925.35 & 21 & 521 & 253 & 2925.24 & 0.11 & \cite{875} \\
$\nu_3+\nu_4$($A_1$) & 3/2 & 2926.91 & 6 & 4 & 2 & 2926.91 & 0.00 & \cite{875} \\
$\nu_3+\nu_4$($F_2$) & 3/2 & 2927.00 & 22 & 1061 & 530 & 2927.00 & 0.00 & \cite{875} \\
$\nu_2+\nu_3$($F_1$) & 3/2 & 3032.87 & 20 & 325 & 186 & 3032.84 & 0.03 & \cite{875} \\
$\nu_2+\nu_3$($F_2$) & 3/2 & 3033.19 & 19 & 537 & 280 & 3033.15 & 0.04 & \cite{875} \\
$\nu_1+\nu_2$($E  $) & 3/2 & 3034.41 & 16 & 86 & 56 & 3034.32 & 0.09 & \cite{875} \\
$2\nu_1$ ($A_1$) & 2 & 4152.98 & 19 & 81 & 35 & 4153.09 & -0.11 & \cite{876} \\
$\nu_1+\nu_3$ ($F_2$) & 2 & 4153.23 & 21 & 1893 & 565 & 4153.09 & 0.14 & \cite{876} \\
\bottomrule
\end{tabular}
\begin{flushleft}\footnotesize
Звездочкой отмечены значения центров <<темных>> полос, оцененные по анализу спектров высокого разрешения; они не использовались в процедуре подгонки. Здесь $N$ --- полиадное число, $N_{\textrm{пер}}$ --- число отнесенных переходов, $N_{\textrm{ур}}$ --- число верхних энергетических уровней.
\end{flushleft}
\end{table}
